\chapter{Introduction}

In real world systems, such as social networks, biological
networks, transportation networks etc. , some nodes
play a very important role in the propagation of
information, influence, or disease \cite{hafiene2020, chen2012,
farooq2018}. Identifying these \textit{influential nodes} is key
challenge in the network science and has significant practical applications
\cite{hafiene2020}. For instance, in viral marketing, targeting a small
set of highly influential users can lead to large scale information
spread \cite{hafiene2020, guo2026}, in epidemiology, immunizing key
individuals can control the outbreak of a disease \cite{chen2012}.

Traditional methods for finding the influential node depends on hand crafted
centrality measures such as degree centrality,betweenness centrality \cite{farooq2018}, closeness centrality
\cite{farooq2018}, and $k$-shell decomposition. While
computationally efficient, these methods often fail to capture complex
structural patterns, especially in weighted or heterogeneous graphs
\cite{mukhtar, chen2012}. Simulation-based approaches such as the
Susceptible-Infected-Recovered (SIR) model provide
more accurate influence estimates, but are computationally expensive and
do not scale to large networks \cite{guo2026, patel2025}.

Recent advances in graph neural networks (GNNs) and convolutional
learning on graphs \cite{zhang2019gcn} open new possibilities
for learning node representations that encode rich structural information
\cite{tu2025, patel2025}. This project proposes a novel framework, the
\textbf{Neural Local Graph Convolutional Network (NLGCN)}, which
combines structural feature engineering with deep learning to predict
influential nodes in weighted graphs efficiently and accurately
\cite{guo2026}.

\section{Problem Statement}

Given a weighted, undirected graph $G = (V, E)$, the goal is to assign
an influence score to each node $v \in V$ such that nodes with higher
scores are more likely to trigger large-scale information diffusion when
selected as initial spreaders \cite{hafiene2020, chen2012}. The
challenge lies in capturing both local neighbourhood structure and
global network topology \cite{mukhtar, guo2026} without resorting to
expensive simulation procedures at inference time \cite{patel2025}.

\section{Objectives}

The key objectives of this project are as follows:
\begin{itemize}
    \item To design structural feature descriptors --- Node Local
    Influence (NLI) and Node Global Influence (NGI), that capture
    multi-scale structural properties of nodes in weighted graphs
    \cite{guo2026, mukhtar}.

    \item To construct a multi-channel tensor representation for each
    node by embedding these features into neighbourhood subgraph
    matrices at three aggregation orders \cite{guo2026, tu2025}.

    \item To design and train an attention-augmented convolutional neural
    network (NLGCN) that predicts node influence scores from the tensor
    representation \cite{ zhang2019gcn}.

    \item To generate ground truth influence labels using the SIR
    epidemic model and evaluate the proposed model
    using Kendall Tau correlation and Top-$N$ spread
    metrics \cite{guo2026}.
\end{itemize}

\section{Motivation}

The motivation for this work stems from two key observations. First,
existing centrality-based methods are inadequate for capturing the
nuanced influence dynamics in weighted and complex networks
\cite{chen2012, mukhtar, farooq2018}. Second, while GNN-based
approaches have shown promise \cite{zhang2019gcn, patel2025},
they typically require attribute-rich graphs or focus on classification
tasks rather than influence ranking \cite{thotakura2024, tu2025}. The proposed NLGCN framework solves this issue by combining domain-specific
structural features with the representational power of convolutional
learning \cite{guo2026}, acquiring a fast, scalable, and accurate
influential node detection.

\section{Organization of the Report}

The report is organized as follows. Chapter~2
presents a review of existing approaches for influential node detection, including traditional centrality-based methods, deep learning techniques, and hybrid approaches. Chapter~3
describes the proposed methodology in detail, covering graph
representation, feature construction, model architecture, label
generation, training procedure, and evaluation metrics. Chapter~4
shows the experimental results and comparative analysis. Chapter~5
discusses the findings and their implications. Chapter~6 concludes the
report and outlines directions for future work.